problem:  
Let \( (M^{2n}, g) \) be a compact, simply connected Riemannian manifold with **nonnegative sectional curvature** supporting an effective, isometric, **slice‑maximal torus action** by \(T^k\).  
Let \(P = M/T^k\) denote the orbit space endowed with its intrinsic **Alexandrov metric** \(d_g\).  
For many known examples (e.g., quasitoric manifolds, torus orbifolds), \(P\) is a **simple convex polytope** whose facets \(F_i\) correspond to codimension‑two strata, and the isotropy representation of \(T^k\) on a normal slice at a point in \(F_i\) determines an integer primitive vector
\[
\lambda(F_i) \in \mathbb{Z}^k,
\]
collectively forming the **characteristic function** \(\lambda\) in the classical Davis–Januszkiewicz classification.

Known rigidity statements show that the **combinatorial type of \(P\)** and the map \(\lambda\) determine the equivariant diffeomorphism type of \(M\).  
However, it is unknown whether the **metric geometry** of \(P\) alone determines \(\lambda\).

**Open problem (Alexandrov metric rigidity of torus actions):**  
Assume \(P_1 = M_1/T^k\) and \(P_2 = M_2/T^k\) are simple polytopes equipped with their quotient Alexandrov metrics \(d_{g_1}\), \(d_{g_2}\).  
Suppose there exists an **isometry of metric polytopes**
\[
\Phi: (P_1,d_{g_1}) \to (P_2,d_{g_2})
\]
that preserves the face stratification.  

1. Does \(\Phi\) necessarily induce a correspondence between facets determining the same isotropy weight data up to a global transformation \(A \in \mathrm{GL}(k,\mathbb{Z})\); that is,
   \[
   \lambda_2(\Phi(F_i)) = A\,\lambda_1(F_i) \quad \forall\, \text{facets } F_i?
   \]
2. Equivalently, can the **Alexandrov metric and face‑intersection angles** of \(P\) alone recover the lattice labels \(\lambda(F_i)\)?  
If not, can one construct examples of distinct torus actions (non‑equivariantly diffeomorphic) whose orbit spaces are **isometric as Alexandrov metric polytopes**?

Why is it a "good" problem:  
This problem isolates a concrete and mathematically meaningful unknown: whether the **quantitative orbit geometry** in the quotient (embodied in the Alexandrov metric structure) determines the **integral isotropy data** that fully encodes the torus action.  

- **Bridges distinct subfields:** It unifies Alexandrov geometry (metric theory of lower‑curvature‑bounded spaces), transformation group topology (classification of torus actions via characteristic functions), and Riemannian geometry (curvature‑related rigidity).  
- **Precisely framed novelty:** Known classification theorems rely on combinatorial data \((P, \lambda)\); this conjecture asks whether the metric quotient provides enough information to recover \(\lambda\), a question unexplored in the standard frameworks.  
- **Rigor and potential depth:**  
  - A positive solution would unveil a new **metric‑to‑integral rigidity principle**, suggesting that curvature and quotient geometry dictate discrete topological invariants.  
  - A negative solution (explicit isospectral–like counterexamples) would reveal unexpected **flexibility**—different torus actions producing identical quotient geometries—analogous to non‑isometric but isospectral phenomena in spectral geometry.  
- **Clarity and cross‑disciplinary reach:** The conjecture is simple to state but demands sophisticated synthesis of ideas from metric analysis, Lie group actions, and toric topology.

By asking if the smooth, equivariant data of a curved manifold is rigidly encoded in its Alexandrov quotient metric—a question neither settled nor trivial—this formulation meets the criteria of a “good’’ research‑level mathematical problem: concise, deep, cross‑field, and conceptually fertile.